\section{Previous Work}\label{sec:PreviousWork}
We review perspective-correct interpolation, rasterization, hardware without perspective correction, and texture filtering.

Texture mapping goes back to Catmull~\cite{Catmull74SAC}.
Heckbert and Moreton~\cite{Heckbert91IPT} show that the texture coordinates divided by the depth are affine functions of the screen position, and so is the reciprocal of the depth.
Interpolating both and dividing them per pixel gives perspective-correct texture coordinates.
Blinn~\cite{Blinn92HI} derives the same interpolation and calls it hyperbolic.
We use it as the reference and quantify the error of its affine approximation.

Pineda~\cite{Pineda88PAP} rasterizes triangles with edge functions, which also give the screen-space barycentric coordinates.
Olano and Greer~\cite{Olano97TSC} rasterize in 2D homogeneous coordinates and interpolate perspective-correctly without dividing the vertices by their depth.
Our software rasterizer uses edge functions.

The first PlayStation lacks perspective correction.
Some of its games split large polygons into smaller ones to reduce the distortion, others replace textures by solid colors~\cite{Copetti19PAP}.
Our bound tells how fine such a split must be.

Mipmapping~\cite{Williams83PP} filters textures to avoid aliasing.
It is orthogonal to our analysis, which considers the texture coordinates before the lookup.
